% Executive summary and analysis of the borough's response to access request
% 2026-581 (study of rue de Terrebonne). Compile with pdflatex (twice).
\documentclass[11pt,letterpaper]{article}
\usepackage[T1]{fontenc}
\usepackage[utf8]{inputenc}
\usepackage[scaled=0.95]{helvet}
\renewcommand{\familydefault}{\sfdefault}
\usepackage[letterpaper,top=0.9in,bottom=0.95in,left=1in,right=1in]{geometry}
\usepackage{microtype}
\usepackage{enumitem}
\usepackage{titlesec}
\usepackage{fancyhdr}
\usepackage{xcolor}
\usepackage{booktabs}
\usepackage{longtable}
\usepackage{array}
\usepackage[hidelinks]{hyperref}
\urlstyle{same}
\newcommand{\docversion}{2026-10-13-12-26-15} % set by build.py (America/Toronto)
\hypersetup{pdfsubject={version \docversion}}

\setlength{\parindent}{0pt}
\setlength{\parskip}{0.5em}
\linespread{1.03}

\titleformat{\section}{\normalfont\large\bfseries}{\thesection}{0.6em}{}[\vspace{1pt}\titlerule]
\titlespacing*{\section}{0pt}{1.1em}{0.4em}
\titleformat{\subsection}{\normalfont\bfseries}{\thesubsection}{0.6em}{}
\titlespacing*{\subsection}{0pt}{0.8em}{0.25em}

\pagestyle{fancy}
\fancyhf{}
\renewcommand{\headrulewidth}{0pt}
\fancyfoot[C]{\footnotesize\color{gray}Access request 2026-581 -- executive summary and analysis -- page \thepage}
\fancyfoot[R]{\tiny\color{gray}v.\,\docversion}

% ---- Repository links: fill in once the documents are uploaded ----
\newcommand{\repourl}{https://github.com/OpenMobilityData/Terrebonne-access-requests}
\newcommand{\repo}[1]{\href{\repourl}{#1}}
% \doc{path-in-repo}{label}: link to a specific file in the repository
\newcommand{\doc}[2]{\href{\repourl/blob/main/#1}{#2}}
% GitHub Pages URL (Settings > Pages > deploy from main, root). Page anchors
% (#page=N) work there because Pages serves PDFs with the right content type;
% GitHub's blob viewer and raw.githubusercontent.com do not honour them.
\newcommand{\pagesurl}{https://openmobilitydata.github.io/Terrebonne-access-requests}
\expandafter\def\csname bundle01\endcsname{2026-581/raw/01_emails_mandate_counts.pdf}
\expandafter\def\csname bundle02\endcsname{2026-581/raw/02_framework_agreement_tender.pdf}
% \pg{bundle}{pages}: page reference, linked to the bundle at its first page
% \pgl{bundle}{anchor page}{displayed pages}: same with an explicit anchor
\newcommand{\pgl}[3]{\href{\pagesurl/\csname bundle#1\endcsname\#page=#2}{\textsf{\small[#1, p.~#3]}}}
\newcommand{\pg}[2]{\pgl{#1}{#2}{#2}}
\newcommand{\fillin}[1]{\textcolor{red}{[#1]}}
\newcommand{\quot}[1]{\begin{quote}\small #1\end{quote}}

\begin{document}

\begin{center}
{\Large\bfseries Access request 2026-581: what the borough's file shows}\\[0.4em]
{\large Executive summary and analysis of the documents released October~9, 2026}\\[0.6em]
Richard Hoge, C\^ote-des-Neiges--Notre-Dame-de-Gr\^ace -- October 9, 2026 (revised October 11)\\
Published with the source documents at \repo{\repourl}
\end{center}

\section*{Summary}

On September~11, 2026 I asked the borough of C\^ote-des-Neiges--Notre-Dame-de-Gr\^ace
for the records behind the engineering study it had commissioned on rue de
Terrebonne: who asked for it, why, what it was told to examine, what it cost,
and what it delivered. The borough answered on October~9 with 508 pages
(file 2026-581). The study itself is withheld. What was released is enough
to establish the following.

\begin{enumerate}[nosep]
  \item \begin{minipage}[t]{\linewidth}\textbf{The decision came before the study.} On June~10, 2026 the
    borough's engineering division wrote to the City's street-design
    division that the elected officials wished to revisit the 2024 layout
    ``avec l'objectif de redonner la bande de stationnement qui avait \'et\'e
    retranch\'ee,'' that the borough's engineers had \emph{already} issued a
    technical opinion on the two concepts the councillors had sketched, and
    that the councillors wanted ``un avis officiel, neutre et externe.'' The
    mandate eventually given to EXP on August~20 is titled, in part,
    ``\'evaluer la place du stationnement sur rue, dans le but d'en augmenter
    le nombre.''\end{minipage}\vspace{0.35em}
\item \begin{minipage}[t]{\linewidth}\textbf{The City's own cycling unit warned about safety.} The
    Service de l'urbanisme et de la mobilit\'e returned an express analysis on
    June~17; the borough's own summary of it, written five days later, says
    the options ``pr\'esentent plusieurs compromis importants sur le plan de
    la s\'ecurit\'e, du confort des usagers vuln\'erables et de la coh\'erence
    du r\'eseau cyclable.'' That analysis is withheld. The fire department's
    opinion of July~13 is redacted in its entirety.\end{minipage}\vspace{0.35em}
\item \begin{minipage}[t]{\linewidth}\textbf{The winter measure is described as transitional; the
    mandate from the mayor's office was permanent.} On August~13 the borough
    mayor's chief of staff set the mandate in writing: a bidirectional lane
    on one side of Terrebonne, from Loyola to C\^ote-Saint-Luc, ``qui
    respecte les normes minimales pr\'evues par la loi,'' and ``le retour du
    stationnement d'un c\^ot\'e de la rue.'' No winter measure appears. The
    service request a week later added Volet~1, removal of cycling
    facilities in winter, as ``mesure transitoire pour l'hiver
    2026--2027,'' ahead of Volet~2, the permanent bidirectional lane. The
    winter conversion announced on October~6 is the first step of a
    year-round reduction of cycling space that has not been disclosed to
    the public.\end{minipage}\vspace{0.35em}
\item \begin{minipage}[t]{\linewidth}\textbf{The ``study'' is a six-week, \$14{,}755 technical note with a
    council date fixed in advance.} Service request August~20; offer
    August~25; purchase order August~26; kick-off September~3; draft
    September~17; final due October~2; announcement October~6. Borough staff
    told the consultant on August~18 that ``l'objectif de nos \'elus est de
    pouvoir agir sur vos conclusions pour un horizon d\'ebut octobre'' and on
    September~10 that the draft was needed by the 18th ``pour respecter notre
    calendrier du Conseil d'arrondissement du 5 octobre.''\end{minipage}\vspace{0.35em}
\item \begin{minipage}[t]{\linewidth}\textbf{The bicycle counts cited at council rest on the weakest
    data in the file and understate winter use; the file contained
    continuous counter data for last winter, documenting average winter usage
    at least five times higher than the cited figure. The more robust data
    were disregarded in the councillor's announcement.} The ``12'' and ``8'' come from the borough's February~2025
    count presentation: not bicycles per day or per hour, but bicycles
    counted using a temporary camera (Miovision) installed at the Cavendish
    and Grand intersections over six
    peak-period hours of a single day, January~29, 2025, in the lane's first
    winter, when clearing was often delayed, and \emph{a day on which Environment
    Canada had issued a winter weather advisory and residents were urged to
    postpone non-essential travel}; the same camera recorded cars on the
    Terrebonne approaches and pedestrians down by about a quarter from the
    November count, bicycles by 95\%. The slide states none of this. Before the first email in the file (June~10) and the
    mandate from the mayor's office (August~13), a continuous counter had
    already recorded every bicycle on the street through the whole of the
    following winter, a full year more recent; the borough's engineers
    offered its data to the City on June~10, and its table was in the
    consultant's hands in August: 4{,}891 passages over December--March
    (41 a day; 46 on weekdays) and a second summer at about 700 a day. For
    the question that bears on safety, exposure, the relevant figure is the
    total number of winter trips, not the count on one chosen day; that
    total was in the borough's possession and was not cited at council. The ``99\%'' drop cited at the meeting, 700
    against 8, sets the largest number on one slide, Grand in
    September~2024 counted on all four approaches, against the smallest on
    the next, Grand in January~2025 on the Terrebonne axis alone; compared
    like for like, the deck's own slides give 97--98\%, label January
    $-95$\%, and the counter gives about 93\% for the seasons entire. The
    two sources agree that the drop is steep;
    they differ in the absolute number, and only one of them yields ``four
    people.'' Whatever the
    explanation, a decision presented as data-driven rested on six hours of
    counting on a single day and ignored thousands of hours of more recent
    counting over an entire winter.\end{minipage}\vspace{0.35em}
\item \begin{minipage}[t]{\linewidth}\textbf{The need for additional winter parking was never measured.}
    Across the 508 pages there is no inventory of on-street spaces, no
    occupancy count, no survey of residents, and no estimate of the spaces
    the conversion would create; the ``50 to 60 spaces'' the councillor gave
    the APC-NDG at its consultation meeting appears nowhere in the release. The mandate asks the consultant to count
    the spaces an arrangement ``pourrait permettre de recr\'eer'' and to
    measure ``l'am\'elioration de la disponibilit\'e du stationnement''; the
    offer promises an analysis of ``l'impact sur le nombre de places'' and
    lists no field work; the consultant worked from aerial imagery. Whatever
    the note says about parking is a count of spaces that could be drawn,
    not a measurement of whether they are needed. Nothing in the file
    records how many winter nights parking near Terrebonne is actually
    full.\end{minipage}\vspace{0.35em}
\item \begin{minipage}[t]{\linewidth}\textbf{The case for removing a lane rests on a view that the lanes are
    wider than their bicycle traffic justifies; that view is the
    councillors' opinion, not the engineers' ruling.} The June~10 email
    reports the elected officials' concern that the lanes' width was ``jug\'ee
    disproportionn\'ee par rapport au volume de cyclistes observ\'e \ldots\
    selon les \'elus,'' and Councillor Teodoresco repeated the point to the
    APC-NDG at the consultation meeting, though not at the October~6 council
    meeting. No engineer in the file endorses it. The borough's own engineers
    describe the alternatives as ``une voie de circulation tr\`es \'etroite
    (moins de 4 m\`etres)'' and ``un am\'enagement bidirectionnel tr\`es
    \'etroit''; the City's cycling unit found ``plusieurs compromis
    importants''; and the August~13 mandate asks for a design ``qui respecte
    les normes minimales pr\'evues par la loi,'' which concedes that the June
    concepts did not. Lane width is set by the cyclist's operating envelope,
    by passing room and by the snow-clearing equipment, none of which depend
    on how many people ride; nothing in the file describes the existing
    lanes as wider than the standard requires.\end{minipage}\vspace{0.35em}
\item \begin{minipage}[t]{\linewidth}\textbf{A storm-only option was within the mandate.} EXP's offer
    lists among the scenarios to compare ``la fermeture en p\'eriode de
    d\'eneigement seulement.'' The councillor told the October~6 meeting that
    storm-only parking had been ruled out; the analysis on which that rests
    is in the withheld note.\end{minipage}\vspace{0.35em}
\item \begin{minipage}[t]{\linewidth}\textbf{The exercise does not follow the City's own published method
    for reviewing a bike lane.} On February~10, 2026 the administration
    launched its audit of the cycling network on the principle ``recueillir
    les donn\'ees, analyser et corriger au besoin'': technical analysis, data
    collection and field observation through the peak season, led by the
    Service de l'urbanisme et de la mobilit\'e, with results expected at the
    end of 2026; the same release names Terrebonne among streets where the
    City is working with the borough on safety. On Terrebonne the objective
    was fixed before any analysis, the SUM's analysis was unfavourable and
    withheld, the mandate was to increase parking, no data were collected,
    and the change was announced while the City's audit was still under
    way.\end{minipage}\vspace{0.35em}
\item \begin{minipage}[t]{\linewidth}\textbf{The withholding of the technical report rests on a ground
    that expires with the decision.} The note and the City's analysis are withheld under
    section~39, which applies ``until a decision is made on the
    recommendation.'' The councillor announced the decision on October~6 and
    promised changes by Christmas. Either the decision has been made and the
    ground has lapsed, or it has not and the announcement was premature. A
    review may be sought from the Commission d'acc\`es \`a l'information
    within 30 days, that is by \textbf{November~8, 2026}.\end{minipage}\vspace{0.35em}
\end{enumerate}

\section{What was asked and what was received}

The request (September~11, 2026; acknowledged September~14 as file 2026-581)
had five parts: A, origin and purpose; B, procurement and contract; C,
methodology and execution; D, deliverables and review; E, internal handling
and communications. The response letter answers each item in turn and
attaches eight files. The two substantive bundles are:

\begin{itemize}[nosep]
  \item \textbf{Bundle 01} (\doc{2026-581/raw/01_emails_mandate_counts.pdf}{01 Doc\_req
    2026-581.pdf}, 167 pages): the email chains, the service request, EXP's
    offer of services, the purchase order, the kick-off meeting form, the
    February~2025 count-analysis presentation and its 2025 update, the
    Kensington loop daily totals (July~26, 2025 to August~31, 2026) and the
    Miovision intersection count tables for 2024--2025.
  \item \textbf{Bundle 02} (\doc{2026-581/raw/02_framework_agreement_tender.pdf}{02 DOC requis
    2026-581.pdf}, 341 pages): the public tender for the 2025--2027
    framework agreements (59-25-ET-01), the contract-management by-law, the
    decision file and selection-committee results, and WSP's winning bid for
    contract~1. Nothing in it concerns Terrebonne specifically; it documents
    the vehicle under which EXP was retained.
\end{itemize}

The released files are published as received in the \repo{OpenMobilityData/Terrebonne-access-requests} repository on GitHub, under \texttt{2026-581/raw/}; derived material, including this document, is kept apart under \texttt{2026-581/derived/}. The remaining attachments are the council resolution awarding the framework
agreements (CA25 170126, June~9, 2025), the notice of recourse, and copies
of the sections of the Act relied on (9, 23, 24, 37, 39, 53--59). The
response letter is at \doc{2026-581/raw/response_2026-10-09.pdf}{Request
2026-581}; the original request is at
\doc{2026-581/raw/request_2026-09-11.pdf}{Terrebonne study access request}.


\textbf{Items with no documents.} The borough found nothing for A2 (external
funding), C3 (follow-up meetings), E2 (communication plans) and E3 (design or
implementation documents dated before receipt of the final report). C3 is
hard to square with the released emails, which schedule working meetings
no.~1 and no.~2 on Teams and refer to the consultant's notes of them
\pgl{01}{23}{23--24, 167}; the kick-off form records that ``on pourra produire
des comptes rendus distincts pour les prochaines rencontres.'' Either those
minutes were not produced or they were not searched for.

\section{Timeline established by the record}

\begin{small}
\begin{longtable}{>{\raggedright\arraybackslash}p{2.6cm}>{\raggedright\arraybackslash}p{10.2cm}>{\raggedright\arraybackslash}p{2.0cm}}
\toprule
Date & Event & Source \\
\midrule
\endhead
2022--2023 & EXP traffic study for the protected lane (previous mandate); scenario retained: one-way with diverging direction at Cavendish, one parking lane removed. & \pg{01}{8} \\
2024 & Redesign built; 2.5~km one-way, unidirectional lanes both sides. & \pg{01}{3} \\
Jul 2024 -- Jun 2025 & Borough Miovision counts at Cavendish, Grand, Walkley; Telraam at King Edward and Royal. & \pgl{01}{82}{82--158} \\
Jan 29, 2025 & Single-day Miovision counts at Cavendish and Grand: 18 and 12 bicycles on the Terrebonne axis over 12 hours; 10 and 8 in the two peak periods. & \pgl{01}{85}{85, 155} \\
Feb 2025 & ``Suivi de projet -- Analyse des comptages'' presentation; updated later in 2025 with January and June counts and the $-95$/$-96$\% figures. & \pgl{01}{27}{27--73} \\
Jul 26, 2025 & Kensington loop counter in service. & \pg{01}{74} \\
Nov 2025 & Municipal election. & --- \\
Jun 10, 2026 & Division head (Trottier) to SUM: elected officials want to restore the parking lane; two cross-sections drawn (designated roadway; very narrow bidirectional lane); borough has already issued a technical opinion; councillors want an external one; meeting with the borough mayor set for June~18. & \pgl{01}{7}{7--8} \\
Jun 17, 2026 & SUM returns its express technical analysis (withheld). & \pgl{01}{6}{6--7} \\
Jun 22, 2026 & Borough asks SIM for an opinion; summarizes SUM's conclusion as ``plusieurs compromis importants'' on safety. & \pgl{01}{5}{5--6} \\
Jul 13, 2026 & SIM replies (body redacted). & \pg{01}{5} \\
Aug 13, 2026 & Chief of staff (D'Amours) to the director of urban planning (Admo), cc.\ borough director, borough mayor and Councillor Teodoresco: the mayor wants a formal external study; mandate: bidirectional lane on one side, Loyola to C\^ote-Saint-Luc, meeting ``les normes minimales pr\'evues par la loi,'' and return of parking on one side. Forwarded to the engineering division the same day. & \pgl{01}{1}{1--2} \\
Aug 18, 2026 & Soulière to EXP: the new borough mayor wants a technical note; speed is what matters to the councillors; aim is to act at a council meeting in early October. & \pgl{01}{21}{21--22} \\
Aug 20, 2026 & Service request 59-26-ET-01\_002 to EXP. & \pgl{01}{3}{3--4} \\
Aug 25, 2026 & EXP offer: \$14{,}755; final note October~2. & \pgl{01}{9}{9--13} \\
Aug 26, 2026 & Purchase order 1796945, \$17{,}706 with contingency. & \pgl{01}{15}{15--17} \\
Aug 28, 2026 & Borough sends EXP the 2025 count presentation and the loop-counter table; EXP sends its first interim analysis (south side preferred for a bidirectional lane). & \pg{01}{19} \\
Sep 3, 2026 & Kick-off meeting (Teams). & \pg{01}{26} \\
Sep 10, 2026 & Borough: draft needed by September~18 ``pour respecter notre calendrier du Conseil d'arrondissement du 5 octobre.'' & \pg{01}{23} \\
Sep 11, 2026 & This access request filed. & --- \\
Sep 11 \& 14, 2026 & Working meetings no.~2 (parking and winter) and comments. & \pg{01}{167} \\
Sep 17, 2026 & Preliminary note delivered, watermarked ``BROUILLON.'' & \pg{01}{167} \\
Oct 2, 2026 & Final note due. No transmittal for the final version appears in the release. & \pgl{01}{12}{12, 26} \\
Oct 6, 2026 & Council meeting; winter conversion announced. & --- \\
Oct 8, 2026 & Second access request (final note and slides) filed. & --- \\
Oct 9, 2026 & Response to this request. & --- \\
\bottomrule
\end{longtable}
\end{small}

\section{Findings in detail}

\subsection{The objective was fixed before any analysis}

The June~10 email from the head of the borough's technical studies division
to the City's street-design division is the earliest document in the
release and the most important \pgl{01}{7}{7--8}:

\quot{``Les \'elus de l'arrondissement souhaitent revoir l'am\'enagement cyclable
de la rue Terrebonne r\'ealis\'e en 2024, avec l'objectif de redonner la bande de
stationnement qui avait \'et\'e retranch\'ee. \ldots\ Les pr\'eoccupations
exprim\'ees concernent principalement la largeur des am\'enagements cyclables,
jug\'ee disproportionn\'ee par rapport au volume de cyclistes observ\'e. Selon
les \'elus, il serait possible de r\'eduire l'espace consacr\'e aux cyclistes
afin de r\'etablir la deuxi\`eme bande de stationnement retir\'ee lors de la
reconfiguration de 2024.''}

\quot{``Vous comprendrez que nous avons d\'ej\`a \'emis un avis technique quant
\`a la faisabilit\'e de ces 2 concepts, toutefois, les \'elus souhaitent obtenir
un avis officiel, neutre et externe \`a l'arrondissement.''}

Three things follow. The objective, restoring the parking lane, is stated as
the elected officials' goal, not as a question. The borough's own engineers
had already given an opinion, and the request for an outside one was made
because the councillors wanted another; the sentence does not say what the
internal opinion was, but the email's own description of the two concepts
(``une voie de circulation tr\`es \'etroite (moins de 4 m\`etres) coinc\'ee
entre deux bandes de stationnement''; ``un am\'enagement bidirectionnel
tr\`es \'etroit'' at 2.4~m with a 0.3~m buffer) and its anticipation that the
fire department ``pourrait s'y opposer'' leave little doubt. And the
justification offered for the exercise is the one used in public since:
the lane is too wide for the number of cyclists.

The email ends with a postscript that bears on the figures used later:
``Nous disposons \'egalement d'un compteur permanent \'Eco-Compteur sur le
corridor. Au besoin, nous pourrons vous transmettre les donn\'ees
disponibles'' \pg{01}{8}. In June 2026 the borough's engineers knew the
continuous counter existed, regarded it as the relevant evidence, and
offered its data to the City.

The internal ``avis technique'' was not covered by my request and has not
been asked for. It should be.

The next document in sequence is the mayor's office's instruction of
August~13 \pgl{01}{1}{1--2}, from the chief of staff to the director of
urban planning, copied to the borough director and both elected officials:

\quot{``Apr\`es discussion, Mme la mairesse souhaite qu'une \'etude technique
formelle soit r\'ealis\'ee par une firme externe afin d'\'evaluer la
reconfiguration de la piste cyclable sur la rue Terrebonne. Le mandat serait
le suivant : \'evaluer le sc\'enario d'implantation d'une piste cyclable
bidirectionnelle d'un c\^ot\'e de la rue Terrebonne (entre le campus Loyola
et C\^ote-Saint-Luc) qui respecte les normes minimales pr\'evues par la loi
et le retour du stationnement d'un c\^ot\'e de la rue.''}

Two things stand out. The mandate as the cabinet framed it is permanent and
full-length, a bidirectional lane on one side and parking restored on the
other, with no mention of winter. And ``qui respecte les normes minimales
pr\'evues par la loi'' presupposes that the designs considered in June did
not: the concepts described on June~10 were a traffic lane under 4~m
between two parking rows and a 2.4~m bidirectional lane, the borough's
engineers had already given their opinion on them, and the SUM had found
``plusieurs compromis importants.'' The August instruction asks for a
version that clears the legal minimum. The winter removal is absent from it;
it enters a week later in the service request drafted by staff, as a
transitional measure ahead of the permanent one. In conversation with the
author, Councillor Teodoresco has said that the borough's engineers would
not endorse a bidirectional lane on Terrebonne, legal though it might be;
the record is consistent with that, and the question it leaves is whether
Volet~2 remains the plan.

\subsection{The City's cycling unit and the fire department}

The SUM's analysis of June~17 \pgl{01}{6}{6--7} is withheld under sections~37
and 39. Its conclusion survives in the borough's own words, in the June~22
request to the fire department \pg{01}{6}:

\quot{``\`A la demande de l'arrondissement, le Service de l'urbanisme et de la
mobilit\'e a proc\'ed\'e \`a une analyse technique pr\'eliminaire des sc\'enarios
envisag\'es. Cette analyse conclut que les options \'etudi\'ees pr\'esentent
plusieurs compromis importants sur le plan de la s\'ecurit\'e, du confort des
usagers vuln\'erables et de la coh\'erence du r\'eseau cyclable par rapport \`a
l'am\'enagement actuellement en place.''}

The SUM's cover note asks to be kept informed ``des orientations qui seront
retenues \`a la lumi\`ere de ces \'el\'ements'' \pg{01}{7}; the release contains
no reply. The fire department's answer of July~13 \pg{01}{5} is a single
introductory sentence followed by a redaction covering the rest of the page.
Whatever it says, the borough asked the question because it expected an
objection to a sub-4-metre lane between two rows of parked cars.

The sequence matters: the City's expert unit flagged safety on June~17; the
borough then retained an outside consultant on August~20 with a mandate
whose stated purpose was to increase parking. A note written to that mandate
is not an independent assessment of whether to remove the lane; it is a
design exercise in how.

\subsection{The mandate and what it asked for}

The service request of August~20 \pgl{01}{3}{3--4} is explicit about purpose and
scope. The title: a technical note to evaluate a bidirectional lane on one
side ``et voir \`a \'evaluer la place du stationnement sur rue, dans le but
d'en augmenter le nombre.'' Among the tasks:

\quot{``\'Enum\'erer les mesures \`a mettre en oeuvre, en deux volets, pour le
retrait des am\'enagements cyclables durant la saison hivernale.
Volet~1 -- existant : pistes cyclables unidirectionnelles bilat\'erales
(mesure transitoire pour l'hiver 2026--2027). Volet~1 alternatif : maintien
et entretien hivernal d'une des deux bandes cyclables, et retrait de
l'autre. Volet~2 -- faisabilit\'e d'une piste cyclable bidirectionnelle
unilat\'erale.''}

``Volet~1 alternatif'' is the measure announced on October~6, word for
word. It appears in the borough's instructions to the consultant six weeks
before the consultant reported. And it is labelled transitional: the
longer-term plan in the same document is Volet~2, a single bidirectional
lane with a parking lane restored on the other side year-round. The
council has not been told that the winter measure is the first stage of a
permanent reconfiguration, but the mandate says so.

EXP's offer \pgl{01}{9}{9--13} restates the scope and adds the winter scenarios
it would compare:

\quot{``\ldots\ la fermeture compl\`ete de la piste et son remplacement par du
stationnement en hiver; la fermeture en p\'eriode de d\'eneigement seulement;
ou, dans le cas des bandes unidirectionnelles actuelles, la fermeture
partielle avec une chauss\'ee d\'esign\'ee et du stationnement des deux c\^ot\'es
de fa\c{c}on permanente ou seulement lors des p\'eriodes de d\'eneigement
durant la saison hivernale.''}

So the storm-only arrangement the councillor dismissed at the meeting was one
of the scenarios the consultant was paid to assess. Whether the note
recommends against it, and on what grounds, is exactly what the withheld
document would show.

Price: \$14{,}755 in fees plus \$2{,}951 contingency, \$17{,}706 before tax
\pg{01}{16}, at the hourly rates of the framework agreement. Deliverables:
two interim emails, a draft note, a final note, and a closing meeting. No
field work, no counts, no occupancy survey; the consultant was given the
borough's 2025 presentation and the counter table on August~28 and worked
from those and from Google aerial imagery, Street View being out of date
\pg{01}{19}.

\subsection{A calendar set by the council date}

The August~18 email that opens the file \pgl{01}{21}{21--22}:

\quot{``\ldots\ la nouvelle mairesse souhaite que nous lui pr\'esentions une note
technique concernant la faisabilit\'e de r\'ealiser une piste cyclable
bidirectionnelle, unilat\'erale, sur la rue de Terrebonne et voir, en plus de
cette faisabilit\'e, dans quelle mesure cet am\'enagement peut augmenter le
nombre de cases de stationnement sur rue. \ldots\ sachez que l'important
pour nos \'elus est que votre proposition puisse se faire tr\`es rapidement,
ainsi que la r\'ealisation de la note en soi. L'objectif de nos \'elus est de
pouvoir agir sur vos conclusions pour un horizon d\'ebut octobre, lors d'un
conseil d'arrondissement.''}

And on September~10 \pg{01}{23}: ``Mon gestionnaire m'a rappel\'e hier qu'il
est tr\`es important que la note technique pr\'eliminaire nous soit achemin\'ee
au plus tard le 18 septembre prochain pour respecter notre calendrier du
Conseil d'arrondissement du 5 octobre prochain.'' The consultant's own draft
of September~17 carries a ``BROUILLON'' watermark because ``le d\'elai serr\'e
ne nous a pas permis de faire faire une relecture attentive'' \pg{01}{167}.
The final note was due October~2, a Friday; the announcement was made the
following Tuesday.

\subsection{Where the numbers came from}

The councillor said at the meeting that he was relying on ``the official
data provided to me by the services'' and read from slides a figure of 12
trips at peak hour in January and a figure of 8. The slides are in the
release \pgl{01}{51}{51--73}: the count-analysis presentation of February~2025,
updated during 2025. Its bicycle table (``MAJ 2025'') gives, for the two
three-hour peak periods combined, 12 bicycles at Cavendish and 12 at Grand
in all directions in January~2025, and 10 and 8 on the Terrebonne axis,
against 347 and 425 in June~2025, labelled $-95$\% and $-96$\%
\pgl{01}{61}{61--62}. The underlying Miovision tables show the January figures
are from one day, January~29, 2025 \pgl{01}{85}{85, 155}; the full twelve-hour
totals that day were 18 at Cavendish and 12 at Grand on the Terrebonne axis.

Three observations. First, these are single-day camera counts at two
intersections in the lane's first winter, taken before the continuous
counter existed; the slide's heading, \emph{aux heures de pointes}, means
the eight is the total over the two three-hour peak periods (06:30--09:30
and 15:30--18:30 \pg{01}{57}), so it is neither a daily nor an hourly
figure, and the slide itself gives no indication that it is one day's
count. Second, the borough had the continuous counter's record
when the decision was made: the same bundle contains daily totals from
July~26, 2025 to August~31, 2026 \pgl{01}{74}{74--81}; the borough's
engineers had offered the counter's data to the City in June \pg{01}{8}
and sent the table to the consultant on August~28 \pg{01}{19}. It shows
4{,}891 passages between December~1 and March~31 (41 a day; 46 on
weekdays; 63 on weekdays over the November~16 to March~31 clearing season)
and a second summer at about 700 a day. Third, the ``700'' and the ``8''
the councillor set against each other for a ``99\%'' drop are not
comparable figures. The 8 is Grand, Terrebonne axis only, six peak hours,
January~29, 2025 \pg{01}{62}. The only number near 700 on the slides is
707: Grand, September~19, 2024, six peak hours, but \emph{all four
approaches}, including bicycles on Grand itself, which the 8 excludes
\pg{01}{61}; it is also the largest single-intersection count in the
presentation. Dividing the smallest number on one slide by the largest on
the other gives 98.9\%. Compared like for like on the deck's own slides,
same intersection, same axis, same six hours, the drop from September~2024
to January~2025 is 98.2\% (452 to 8) and from the summer counts on either
side of that winter about 97\% (283 and 281 to 8); the slides themselves
label January $-95$\%. The continuous counter gives about 93\% between the
whole of summer and the whole of winter. The drop is steep by every
measure; ``99\%'' is the one pairing of the available numbers that reaches
it.

The day itself. January~29, 2025 was a Wednesday on which Environment and
Climate Change Canada had issued a winter weather advisory for Montr\'eal:
five to ten centimetres of snow with blowing snow, the worst conditions
expected to last until midday, and residents urged to postpone non-essential
travel.\footnote{Jessica Barile, ``Montrealers urged to delay non-essential
travel as snowstorm hits, up to 10 cm possible,'' CityNews Montreal,
January~29, 2025, 9:13~a.m.,
\url{https://montreal.citynews.ca/2025/01/29/montrealers-urged-to-delay-non-essential-travel-as-snowstorm-hits-up-to-10-cm-possible/}.}
The Miovision tables for Grand that day record every mode, not only
bicycles \pgl{01}{82}{82--88}. Over the same six peak hours the camera
counted 2{,}298 cars on the four approaches (617 on the two Terrebonne
approaches), 108 buses, 20 trucks and 605 pedestrians, against the 8
bicycles on the Terrebonne axis, three of them in the morning peak and five
in the afternoon. Set against the previous count at the same intersection,
November~20, 2024 \pgl{01}{131}{131--137}, cars on the Terrebonne
approaches were down 23\% (802 to 617), pedestrians down 26\% (818 to 605)
and bicycles on the axis down 95\% (154 to 8); against September~19, 2024,
cars were down 11\%, pedestrians 31\% and bicycles 98\%. Every mode fell
that day; cycling fell furthest, as it does on any snow day, which is the
reason a single day is not a measure of a season. The slide from which the
8 was read gives neither the date nor the weather. The counts by mode for
every Miovision day in the file are tabulated in
\doc{2026-581/derived/miovision_counts_by_mode_2026-581.xlsx}{derived/miovision\_counts\_by\_mode\_2026-581.xlsx}.

For completeness: the presentation contains no Kensington data, nor could
it, the loop having entered service on July~26, 2025, after the deck's last
update (June~2025). Its methods slide does list the Eco-Counter, ``utilis\'e
par le SUM depuis 2013'', as one of the three count sources alongside
Miovision and Telraam \pgl{01}{53}{29, 53}; the deck's only continuous data
are the citizen Telraam sensors on King-Edward and Royal, which give about
242 bicycles a day \pg{01}{65}. The counter table reached the consultant
as a separate attachment to the same August~28 email as the presentation
\pg{01}{19}. The consultant therefore had both documents from the first
week of the mandate, and the borough had known since at least June which
was the better source \pg{01}{8}; which of them the note relies on cannot
be known until it is released.

\subsection{What the file contains on parking}

The stated purpose of the exercise is more parking, so it is worth
recording what the release contains on that side. The June~10 email states
the objective and gives one geometric figure, the sub-4-metre traffic lane
between two parking rows \pg{01}{8}; nothing on demand. The service request
asks EXP to ``mesurer \ldots\ le nombre de cases de stationnement que cet
am\'enagement pourrait permettre de recr\'eer'' and to measure
``l'am\'elioration de la disponibilit\'e du stationnement sur de
Terrebonne'' \pgl{01}{3}{3--4}: questions put to the consultant, not
answers. EXP's offer promises ``une analyse de l'impact sur le nombre et
l'accessibilit\'e des espaces de stationnement'' and a deliverable
``incluant l'impact sur le nombre de places de stationnement sur rue''; its
list of services contains no occupancy survey and no count of any kind
\pgl{01}{10}{10--11}. The August~28 email shows the first interim analysis
concerned which side of the street should carry a bidirectional lane, with
the parking analysis to follow, and that the consultant was working from
aerial imagery \pg{01}{19}; the kick-off form schedules ``Analyse de c\^ot\'e
de rue du stationnement et am\'enagement hivernal'' for September~11
\pg{01}{26}. The count presentation contains no parking slide. The
response to item~C2 of the request, which asked for the data collection
behind the study ``including locations, dates, seasons and hours of
counts,'' is ``the attached documents include presentations,'' and the
presentations are bicycle and vehicle counts at intersections.

The release therefore contains no measured parking figure of any kind: no
inventory, no occupancy, no survey, no estimate of spaces created. The
``50 to 60 spaces'' the councillor gave the APC-NDG at the meeting he
later cited at council as his consultation with the Association is not in
it. Whatever the
withheld note says about parking will be a geometric count of spaces that
could be marked, which is what the mandate asked for; nothing in the mandate
or the deliverables would produce a measurement of whether, or on how many
nights, those spaces are needed.

\subsection{What was withheld, and on what grounds}

\begin{itemize}[nosep]
  \item \textbf{Section 39} (a study prepared in connection with a
    recommendation, ``until a decision is made on the recommendation''):
    the EXP note in all versions, the SUM analysis, and the comments
    exchanged on deliverables.
  \item \textbf{Section 37} (recommendations or opinions less than ten years
    old, including a consultant's): the same documents, plus the SIM
    opinion.
  \item \textbf{Section 9, second paragraph} (drafts and preliminary notes
    are not documents held by a public body): the interim deliverables and
    the preliminary note.
  \item \textbf{Sections 23--24} (third-party commercial information):
    portions of the tender documents and EXP's offer; the hour-and-fee
    breakdown attached to the offer is not in the release.
  \item \textbf{Sections 53--59}: personal information, applied to names
    and contact details in emails; in practice staff names are left
    visible.
\end{itemize}

Section~39 is the ground that carries the weight, and it is time-limited by
its own terms. The councillor announced the decision on October~6 and told
the Gazette on October~7 that ``one thing is certain: residents will see
changes by Christmas 2026,'' while also saying the change ``still need[s]
to go through several steps.'' A signed statement posted to his
social-media page in the days after the meeting puts both halves in his own
words: ``The solution we are currently developing with experts is to
maintain the bike lane as is during the summer and remove it on one side
only during the winter season. Several steps remain before this proposal
can be implemented, including consideration by the city's decision-making
bodies, namely the borough council and City Council. One thing is certain:
residents will see changes before Christmas 2026.'' The borough cannot have it both ways: if the
recommendation has been decided, section~39 no longer applies and the note
must be released subject only to section~37, which is discretionary and
does not cover factual content; if it has not been decided, the public
announcement of a decided outcome was premature. Section~14 of the Act
requires that factual content be severed and released where an exception
covers only part of a document; the response does not do this for the note.

The section~9 ground is weak for the September~17 draft. It was a
deliverable under contract, transmitted to the client for review, and the
borough's comments on it shaped the final version. ``Brouillon'' on the
watermark does not make it a personal draft within the meaning of
section~9, and the Commission has said so in comparable cases.

\section{What the record does and does not establish}

\subsection{On the decision process}

The record shows a goal set by the elected officials in the spring; an
internal technical opinion already given by the borough's own engineers,
after which an external one was sought; a safety analysis from the City's
cycling unit received in June and not mentioned in public; a mandate written
in August whose stated purpose was to increase parking, with a council date
already chosen; and a presentation of counts that set aside the borough's
continuous counter data in favour of a one-day count from the previous
winter, and no measurement of parking demand at all. Whether the councillors are right to want more parking is a
separate question; what the documents establish is the sequence in which
the decision and its justification were produced, and that sequence differs
from the account given at the October~6 meeting.

\subsection{On the storm-only alternative}

The borough paid for an assessment that included ``la fermeture en
p\'eriode de d\'eneigement seulement'' and stated publicly that the option was
ruled out because residents could not be informed of it. The reasoning, if
there is more to it, is in the withheld note.

\subsection{On what is planned beyond this winter}

The mayor's office asked on August~13 for a permanent bidirectional lane
with parking restored; ``mesure transitoire pour l'hiver 2026--2027'' and
Volet~2 together describe the winter conversion as a first step toward it. That is
a different decision from the one announced, with a different safety
profile: the SUM's ``compromis importants'' concerned exactly these
cross-sections. It has not been disclosed to the public. The councillor's
statement that the lane will be maintained ``as is during the summer'' is
the first public remark bearing on it; it does not say whether the
bidirectional option has been abandoned, and the question should be put
directly. The same statement names City Council among the bodies that must
consider the proposal, which, if so, places the City's own street-design
service, author of the June~17 analysis, on the path to approval.

\subsection{Against the City's own review method}

\textbf{Two mandates.} The lane itself was built after a different kind of
study. The administration elected in 2021 had campaigned on expanding the
cycling network, and the 2023 EXP study it commissioned on Terrebonne was a
feasibility and design study for a project already decided: its task was
to confirm there was no obstacle and to settle the layout, and its
limitations, which opponents of the lane pointed out at the time, were
those of a design study. The administration elected in 2025 campaigned on
something else: not on removing lanes but on ``a comprehensive audit of
all Montreal cycling infrastructure within the first 100 days,'' to
``review, secure or remove cycling infrastructure where necessary,'' with
the leader describing herself as ``pro-bikes.''\footnote{CBC News, ``How
do Montreal's top mayoral candidates compare on election issues?''
(platform comparison, October~2025),
\url{https://newsinteractives.cbc.ca/features/2025/montreal-municipal-election/};
CBC News, ``Are Montreal bike paths a boon or burden? See what mayoral
hopefuls are promising,'' \url{https://www.cbc.ca/1.7648562}; La Presse,
``Martinez Ferrada veut revoir l'ensemble du r\'eseau cyclable,''
August~11, 2025,
\url{https://www.lapresse.ca/actualites/grand-montreal/2025-08-11/elections-municipales/martinez-ferrada-veut-revoir-l-ensemble-du-reseau-cyclable.php}.}
A study that is to decide whether a lane stays carries a weight that a
design study does not, and the administration set out what such a study
involves on February~10, 2026.

The administration's standard for reviewing a bike lane is on the public
record. The City's press release of February~10, 2026 launching its audit of
the cycling network describes the method as ``recueillir les donn\'ees,
analyser et corriger au besoin,'' with ``analyse technique, collecte de
donn\'ees et observation terrain'' on four existing lanes (Rachel,
C\^ote-Sainte-Catherine, de la Commune, Meilleur) through the peak season,
led by the Service de l'urbanisme et de la mobilit\'e, results expected
toward the end of 2026, and a revision of the Plan v\'elo to set ``des
crit\`eres clairs, coh\'erents et harmonis\'es''; the mayor's stated reason
was that ``les documents qu'on nous a remis ne r\'epondaient pas aux
pr\'eoccupations exprim\'ees sur le terrain,'' so the City would collect its
own data first. The same release lists Terrebonne, with Hochelaga, Gouin and
Henri-Bourassa, among streets where the City is working with the borough on
safety.\footnote{Ville de Montr\'eal, ``10 actions en 100 jours : la Ville de
Montr\'eal annonce un \'etat des lieux pour une mobilit\'e planifi\'ee,''
February~10, 2026,
\url{https://montreal.ca/communiques-de-presse/10-actions-en-100-jours-la-ville-de-montreal-annonce-un-etat-des-lieux-pour-une-mobilite-planifiee};
Journal M\'etro, ``Montr\'eal lance son audit des pistes cyclables,''
February~10, 2026,
\url{https://journalmetro.com/actualites/politique/3208452/montreal-lance-son-audit-des-pistes-cyclables/}.}
Measured against that method, the Terrebonne file shows the objective set
before analysis, the SUM's analysis unfavourable and withheld, a mandate
whose purpose was parking, no data collection of any kind, and an
announcement made while the City's audit was still running. Whether the two
processes are formally connected is a question for the administration; the
contrast in method is a matter of record.

\subsection{What the record does not establish}

The SUM's conclusion is known only as the borough summarized it; the SIM's
opinion is unknown; the content of EXP's note is unknown, and it may well
recommend against the narrower options and for the winter measure as the
least bad. ``Disproportionn\'ee'' is staff's report of the councillors'
view, not a finding. None of the quotations above should be stretched past
what they say. They establish the sequence and the purpose; they do not
establish what the consultant concluded, which is the reason to get the
note.

\section{Document index}

Page numbers refer to the PDF page of each released bundle.

\begin{small}
\begin{longtable}{>{\raggedright\arraybackslash}p{3.0cm}>{\raggedright\arraybackslash}p{9.2cm}>{\raggedright\arraybackslash}p{2.6cm}}
\toprule
Document & Content & Location \\
\midrule
\endhead
Response letter & Item-by-item answers; grounds for refusal; notice of recourse. & \doc{2026-581/raw/response_2026-10-09.pdf}{Request 2026-581}, 8 pp. \\
Bundle 01 & & \doc{2026-581/raw/01_emails_mandate_counts.pdf}{01 Doc\_req}, 167 pp. \\
\quad Mayor's office instruction & Aug 13 email from the chief of staff setting the mandate (bidirectional lane, return of parking), forwarded Aug 13--14. & pp. 1--2 \\
\quad SIM / SUM chain & Jun 10 request to SUM; Jun 12, 17 replies; Jun 22 request to SIM; Jul 13 SIM reply (redacted). & pp. 5--8 \\
\quad Service request & Aug 20, 2026, mandate 59-26-ET-01\_002. & pp. 3--4 \\
\quad EXP offer & Aug 25, 2026; scope, scenarios, schedule, \$14{,}755. & pp. 9--13 \\
\quad Purchase order & Aug 26, 2026, no.~1796945, \$17{,}706. & pp. 15--17 \\
\quad EXP email chain & Aug 18--28: mandate, work plan, inputs, first interim analysis. & pp. 19--22 \\
\quad Kick-off & Sep 3 meeting form, deliverables and dates; Sep 10 emails on the Oct 5 council deadline. & pp. 23--26 \\
\quad Count presentation & ``Suivi de projet -- Analyse des comptages,'' Feb 2025, and ``MAJ 2025'' update. & pp. 27--49; 51--73 \\
\quad Loop counter & Kensington daily totals, bicycles and motor vehicles, Jul 26, 2025 -- Aug 31, 2026. & pp. 74--81 \\
\quad Miovision tables & Intersection counts by 15-minute period and class: Grand, Cavendish, Walkley, Royal; Jul 2024 -- Sep 2025; Jan 29, 2025 at pp. 82--88 and 152--158. & pp. 82--164 \\
\quad Draft delivered & Sep 17, 2026 transmittal of the preliminary note. & p. 167 \\
Bundle 02 & Framework-agreement tender 59-25-ET-01, by-law, decision file (pp. 169--177), evaluation grid, WSP bid. & \doc{2026-581/raw/02_framework_agreement_tender.pdf}{02 DOC requis}, 341 pp. \\
Resolution & CA25 170126, Jun 9, 2025: framework agreements to WSP (\$1.60M) and EXP (\$1.70M). & \doc{2026-581/raw/03_resolution_CA25-170126.pdf}{03 R\'esolution}, 2 pp. \\
Act extracts & ss.~9, 23--24, 37, 39, 53--59; notice of recourse. & \href{\repourl/tree/main/2026-581/raw/act_extracts}{act\_extracts/}, 5 files \\
Derived & Kensington daily totals parsed to CSV from pp. 74--81. & \doc{2026-581/derived/kensington_daily_borough_release_2025-07_2026-08.csv}{CSV} \\
Derived & Miovision counts by mode (cars, buses, trucks, bicycles, pedestrians), six peak hours and 12~h, all approaches and Terrebonne axis, every count day, from pp. 82--164; Jan 29, 2025 at Grand broken out. & \doc{2026-581/derived/miovision_counts_by_mode_2026-581.xlsx}{XLSX}, \doc{2026-581/derived/miovision_counts_by_mode_2026-581.csv}{CSV} \\
\bottomrule
\end{longtable}
\end{small}

\end{document}
